% 04_architecture.tex — Rail Paper: PWM Architecture
% ============================================================

\section{PWM Architecture: The Evaluation Harness}
\label{sec:architecture}

This section presents the concrete architecture of the \pwm{} (PWM)
evaluation harness.
PWM comprises four interlocking subsystems: the OperatorGraph IR for
physics representation (\cref{sec:operatorgraph}), the 4-scenario
evaluation protocol for standardized assessment (\cref{sec:scenarios}),
the LIP-Arena for prospective evaluation (\cref{sec:lip_arena}), and
the Red Team module for adversarial verification (\cref{sec:redteam}).
Together, these subsystems implement all nine layers of the Industrial
Intelligence Stack described in \cref{sec:stack}.

% ==================================================================
%  OperatorGraph IR
% ==================================================================

\subsection{OperatorGraph IR: A Universal DAG for Physics}
\label{sec:operatorgraph}

At the core of PWM lies the \textbf{OperatorGraph Intermediate
Representation (IR)}---a directed acyclic graph (DAG) formalism that
represents any forward measurement operator as a composition of
primitive physical operations.
The design is motivated by a key observation: despite the enormous
diversity of imaging modalities, the underlying physics decomposes
into a surprisingly small set of recurring primitives.

\paragraph{Nodes.}
Each node in an OperatorGraph wraps a single primitive operator drawn
from a validated library.
Primitive operators include:
\begin{itemize}
  \item \textbf{Mask modulation}: element-wise multiplication by a
        coded aperture pattern, $\mathbf{x} \mapsto \mathbf{M} \odot
        \mathbf{x}$.
  \item \textbf{Convolution}: linear shift-invariant filtering,
        $\mathbf{x} \mapsto \mathbf{h} \ast \mathbf{x}$, encompassing
        point spread functions, blur kernels, and diffraction models.
  \item \textbf{Spectral dispersion}: wavelength-dependent spatial
        shift, $\mathbf{x}_\lambda \mapsto \mathcal{S}_\lambda
        (\mathbf{x}_\lambda)$, the core operation in CASSI systems.
  \item \textbf{Temporal modulation}: frame-dependent coding,
        $\mathbf{x}_t \mapsto \mathbf{C}_t \odot \mathbf{x}_t$,
        as in CACTI and other snapshot compressive video systems.
  \item \textbf{Projection}: line-integral or Radon transform,
        $\mathbf{x} \mapsto \mathcal{R}_\theta(\mathbf{x})$, for
        tomographic modalities.
  \item \textbf{Fourier sampling}: $k$-space under-sampling,
        $\mathbf{x} \mapsto \mathbf{P}_\Omega
        \mathcal{F}(\mathbf{x})$, for MRI and related modalities.
  \item \textbf{Noise injection}: additive, multiplicative, or
        Poisson noise models, $\mathbf{y} \mapsto \mathbf{y} +
        \boldsymbol{\eta}$.
  \item \textbf{Sensor integration}: spatial and temporal binning,
        quantization, and detector response, modeling the
        analog-to-digital conversion chain.
\end{itemize}

\paragraph{Edges.}
Edges define data flow between nodes.
An edge from node $u$ to node $v$ means that the output of operator $u$
is an input to operator $v$.
The DAG structure enforces a well-defined execution order and prohibits
cycles, ensuring that every OperatorGraph compiles to a deterministic
forward model.
Where a node requires multiple inputs (\eg, a sensor that integrates
both signal and dark current), multiple incoming edges are supported
with explicit concatenation or summation semantics.

\paragraph{Compilation.}
Given an OperatorGraph $\mathcal{G} = (\mathcal{V}, \mathcal{E})$, the
system compiles the graph into a callable forward model
$\mathbf{H}_\mathcal{G}: \mathcal{X} \to \mathcal{Y}$ by topological
traversal.
When all primitive nodes support an adjoint, the system also compiles
$\mathbf{H}_\mathcal{G}^\top$ automatically.
Automatic differentiation through the graph enables gradient-based
calibration of any mismatch parameter.

\paragraph{Physical carriers.}
The primitive library spans five physical carriers:
\begin{enumerate}
  \item \textbf{Photons}: optical imaging (CASSI~\cite{wagadarikar2008cassi}, SPC~\cite{duarte2008spc}, lensless,
        coded diffraction, \etc{}).
  \item \textbf{Electrons}: electron microscopy (cryo-EM, STEM, SEM).
  \item \textbf{Spins}: magnetic resonance imaging~\cite{lustig2007sparse}
        (MRI, diffusion MRI, spectroscopic MRI).
  \item \textbf{Acoustic waves}: ultrasound, photoacoustic imaging,
        sonar.
  \item \textbf{Particles}: computed tomography~\cite{feldkamp1984practical},
        positron emission tomography (PET),
        single-photon emission CT (SPECT), neutron imaging.
\end{enumerate}

\paragraph{Templates and coverage.}
PWM includes 89 validated OperatorGraph templates covering 64 distinct
modalities.
Some modalities require multiple templates to capture different
hardware configurations (\eg, single-disperser vs.\ dual-disperser
CASSI).
Each template specifies the graph topology, the primitive operators at
each node, the expected parameter ranges, and the physically meaningful
mismatch axes.

\paragraph{Example: CASSI.}
The following equation illustrates the OperatorGraph for a standard
single-disperser CASSI~\cite{wagadarikar2008cassi} system:
\begin{equation}
  \mathcal{G}_{\text{CASSI}}:
  \quad
  \underbrace{\mathbf{x}_\lambda}_{\text{Source}}
  \;\xrightarrow{\text{Mask}}\;
  \underbrace{\mathbf{M} \odot \mathbf{x}_\lambda}_{\text{Modulated}}
  \;\xrightarrow{\text{Dispersion}}\;
  \underbrace{\mathcal{S}_\lambda(\mathbf{M} \odot
  \mathbf{x}_\lambda)}_{\text{Dispersed}}
  \;\xrightarrow{\text{Sensor}}\;
  \underbrace{\textstyle\sum_\lambda
  \mathcal{S}_\lambda(\mathbf{M} \odot
  \mathbf{x}_\lambda)}_{\text{Measurement}}
  \;\xrightarrow{\text{Noise}}\;
  \mathbf{y}
\end{equation}
The mismatch axes for this template include mask shift
($\Delta x$, $\Delta y$), mask rotation ($\Delta\theta$), dispersion
slope ($\Delta s$), dispersion axis offset ($\Delta\phi$), detector
gain ($\Delta g$), and additive noise variance ($\sigma^2$).
Each axis has a physically motivated default range derived from
hardware specifications and optical tolerance analysis.

% ==================================================================
%  4-Scenario Protocol
% ==================================================================

\subsection{The 4-Scenario Evaluation Protocol}
\label{sec:scenarios}

Every reconstruction method evaluated within PWM is assessed under four
standardized scenarios that systematically vary the relationship between
the measurement operator and the reconstruction operator.
\Cref{tab:scenarios} defines all four scenarios.

\begin{table}[t]
\centering
\caption{%
  \textbf{The 4-scenario evaluation protocol.}
  $\mathbf{H}$ denotes the true forward operator,
  $\mathbf{H}_{\mathrm{nom}}$ the nominal (assumed) operator,
  and $\hat{\mathbf{H}}$ the calibrated operator.
  Each scenario isolates a different aspect of reconstruction
  performance.%
}
\label{tab:scenarios}
\smallskip
\renewcommand{\arraystretch}{1.15}
\begin{tabular}{@{}clccp{5.2cm}@{}}
\toprule
  & \textbf{Scenario}
  & \textbf{Measurement}
  & \textbf{Reconstruction}
  & \textbf{Interpretation} \\
\midrule
I   & Ideal      & $\mathbf{H}$    & $\mathbf{H}$
    & Oracle upper bound: what is achievable with perfect operator
      knowledge. \\
II  & Assumed/Mismatch    & $\mathbf{H}$    & $\mathbf{H}_{\mathrm{nom}}$
    & Mismatch baseline: the default deployment condition where the
      nominal operator differs from truth. \\
III & Corrected  & $\mathbf{H}$    & $\hat{\mathbf{H}}$
    & Calibration benefit: how much performance is recovered by
      estimating the true operator. \\
IV  & Oracle Mask & $\mathbf{H}$ & $\mathbf{H}$ (mask only)
    & Partial oracle bound: true mask with nominal dispersion, isolating
      the mask contribution. \\
\bottomrule
\end{tabular}
\end{table}

\paragraph{Scenario I (Ideal).}
The true operator $\mathbf{H}$ is used for both measurement generation
and reconstruction.
This is the standard evaluation setting in virtually all published
work and represents the oracle upper bound on reconstruction quality.
It answers the question: \emph{how good is this solver when the physics
is perfectly known?}

\paragraph{Scenario II (Assumed/Mismatch).}
Measurements are generated with the true operator $\mathbf{H}$, but
reconstruction uses the nominal operator
$\mathbf{H}_{\mathrm{nom}}$---the operator that the system
\emph{believes} is correct but that differs from ground truth due to
alignment errors, manufacturing tolerances, or calibration drift.
This is the realistic deployment condition and is the scenario most
relevant to practitioners.
It answers the question: \emph{how robust is this solver to operator
error?}

\paragraph{Scenario III (Corrected).}
Measurements are generated with $\mathbf{H}$, but reconstruction uses a
calibrated operator $\hat{\mathbf{H}}$ obtained by running a
calibration procedure on the nominal operator.
The difference between Scenario~III and Scenario~II measures the
\emph{calibration gain}---the value added by operator correction.
The difference between Scenario~I and Scenario~III measures the
\emph{residual gap}---the remaining performance lost to imperfect
calibration.
It answers: \emph{how much does calibration help?}

\paragraph{Scenario IV (Oracle Mask).}
A partial oracle condition in which the true mask is provided but
dispersion parameters remain at their nominal values.
This scenario isolates the contribution of mask knowledge versus
dispersion knowledge, enabling diagnostic analysis of which mismatch
axis matters most.
It answers: \emph{which operator component is the bottleneck?}

\paragraph{Derived metrics.}
From the four scenario PSNR values, we derive:
\begin{align}
  \text{Degradation}     &= \text{PSNR}_\mathrm{I} - \text{PSNR}_\mathrm{II}, \label{eq:degradation} \\
  \text{Calibration Gain}&= \text{PSNR}_\mathrm{III} - \text{PSNR}_\mathrm{II}, \label{eq:calib_gain} \\
  \text{Recovery Ratio}  &= \frac{\text{PSNR}_\mathrm{III} - \text{PSNR}_\mathrm{II}}{\text{PSNR}_\mathrm{I} - \text{PSNR}_\mathrm{II}}, \label{eq:recovery} \\
  \text{Oracle Gap}      &= \text{PSNR}_\mathrm{I} - \text{PSNR}_\mathrm{III}. \label{eq:oracle_gap}
\end{align}
The recovery ratio and oracle gap are the primary pass/fail metrics
referenced in Layer~1 of the Industrial Intelligence Stack
(\cref{tab:stack_mapping}): a recovery ratio $\rho$ (as defined in Eq.~\ref{eq:recovery}) of at least 0.80, meaning that calibration recovers at least 80\% of the mismatch-induced degradation, and an oracle
gap $\leq 2$\,dB indicate that calibration is sufficient for
deployment.

% ==================================================================
%  LIP-Arena
% ==================================================================

\subsection{LIP-Arena: Prospective Evaluation via Commit-Measure-Score}
\label{sec:lip_arena}

Retrospective benchmarks---including InverseNet---suffer from an
inherent limitation: the test data exists before the methods are
developed.
No matter how carefully the benchmark is designed, determined
participants can eventually overfit to it through hyperparameter tuning,
architecture search, or implicit memorization.
To address this, PWM includes the \textbf{LIP-Arena} (Leaderboard for Imaging Physics Arena), a prospective evaluation competition in which test
data is generated \emph{after} all submissions are finalized.
The key principle is: \textbf{data created after deadline}---no
memorization possible.

The four-phase Commit-Measure-Score protocol proceeds as follows:

\paragraph{Phase 1: Commit (2 weeks).}
Participating teams submit sealed containers encapsulating their
reconstruction method, along with a declared compute budget (FLOPs,
wall-clock time, and memory).
Submissions are cryptographically hashed and timestamped.
No modifications are permitted after the commit deadline.

\paragraph{Phase 2: Measure (2 weeks).}
New measurement data is generated \emph{after} the commit deadline
using OperatorGraph templates with freshly sampled mismatch parameters.
Test scenes are drawn from a held-out pool that has never been
published.
The measurement operator, mismatch parameters, and test scenes are
sealed until Phase~4.

\paragraph{Phase 3: Execute (1 week).}
All submitted containers are run in an identical, sandboxed compute
environment.
Each container receives only the raw measurements and the nominal
operator $\mathbf{H}_{\mathrm{nom}}$---never the true operator or the
ground-truth images.
Execution is monitored for compute budget compliance:
containers that exceed their declared budget are flagged.

\paragraph{Phase 4: Score (1 week).}
Reconstructions are scored automatically against ground-truth images
using PSNR, SSIM, LPIPS, and spectral-angle mapper (SAM) where
applicable.
All four scenarios are evaluated.
Results, RunBundles, and failure analyses are published in an open
round report.

\subsubsection{Evaluation Tracks}
\label{sec:tracks}

LIP-Arena organizes competition into four tracks that target different
aspects of the reconstruction problem:

\begin{enumerate}
  \item \textbf{Correct}: Given mismatched measurements, produce the
        best reconstruction.
        Methods may use any calibration strategy.
        Primary metric: corrected PSNR (Scenario~III).
  \item \textbf{Diagnose}: Given mismatched measurements, identify
        which mismatch parameters are present and estimate their
        values.
        Primary metric: parameter estimation error.
  \item \textbf{No-GT}: Evaluate reconstruction quality without
        access to ground-truth images.
        Methods must provide calibrated uncertainty estimates.
        Primary metric: correlation between predicted and actual PSNR.
  \item \textbf{Design}: Given a modality and a target recovery ratio,
        design the optimal coded aperture pattern or coding scheme.
        Primary metric: recovery ratio achieved by a reference solver.
\end{enumerate}

\subsubsection{Anti-Goodhart Scoring}
\label{sec:anti_goodhart}

Goodhart's Law---``when a measure becomes a target, it ceases to be a
good measure''---is the central threat to any benchmark.
LIP-Arena mitigates this through three mechanisms:

\begin{enumerate}
  \item \textbf{Prospective dominance weighting.}
        Final scores are computed as $0.7 \times
        \text{PSNR}_{\mathrm{prospective}} + 0.3 \times
        \text{PSNR}_{\mathrm{retrospective}}$, ensuring that
        performance on unseen data dominates the ranking.

  \item \textbf{Gaming penalties.}
        Methods that exhibit statistically significant performance
        differences between prospective and retrospective data (\ie,
        methods that have likely overfit to the retrospective set)
        receive a scoring penalty proportional to the gap.

  \item \textbf{Multi-metric ranking.}
        Rankings are computed over an ensemble of metrics (PSNR, SSIM,
        LPIPS, SAM), preventing optimization for any single number.
\end{enumerate}

% ==================================================================
%  Red Team
% ==================================================================

\subsection{Red Team Module: Adversarial Verification}
\label{sec:redteam}

The Red Team module is the adversarial verification layer that
probes submitted methods for failure modes beyond the standard
evaluation scenarios.
It comprises six categories of adversarial tests, totaling over 2{,}900
pre-designed scenarios:

\begin{enumerate}
  \item \textbf{Novel mismatch}: mismatch types not represented in the
        training distribution (\eg, non-rigid mask deformation,
        spatially varying dispersion).
  \item \textbf{Compound mismatch}: simultaneous perturbation of
        multiple operator parameters beyond the range of single-axis
        evaluations.
  \item \textbf{Out-of-family physics}: test scenes or measurement
        conditions drawn from a different physical regime (\eg,
        fluorescence emission applied to an absorption-trained solver).
  \item \textbf{Distribution shift}: test scenes with statistics
        dramatically different from training data (\eg, medical
        images for a solver trained on natural scenes).
  \item \textbf{Compute traps}: inputs designed to trigger worst-case
        computational complexity (\eg, iterative methods that fail to
        converge, attention mechanisms with adversarial token
        distributions).
  \item \textbf{Gate-flip scenarios}: edge cases where a small change
        in the input causes the output to qualitatively change (\eg,
        a mismatch parameter crossing a phase-transition boundary).
\end{enumerate}

Red Team results are reported separately from main evaluation
scores to avoid penalizing methods for failing on out-of-distribution
scenarios.
However, persistent Red Team failures across rounds are flagged in
governance reports and inform the design of future standard
evaluation scenarios---ensuring that the benchmark evolves to cover
newly discovered failure modes.

% ==================================================================
%  Safety Brakes
% ==================================================================

\subsection{Safety Brakes and Governance}
\label{sec:safety}

To prevent the publication of misleading results and ensure that
LIP-Arena maintains scientific integrity, PWM implements pre-committed
safety brakes:

\begin{itemize}
  \item \textbf{Recovery ratio $< 0.30$}: Automatic disqualification.
        A method that recovers less than 30\% of the oracle performance
        after calibration is either fundamentally broken or has not
        been properly configured.  Note: this threshold applies to
        arena submissions, not to inherent modality limitations; some
        modalities (e.g., CASSI) may exhibit recovery ratios below 0.30
        due to the severity of the mismatch, which is diagnosed separately.
  \item \textbf{Uncertainty miscalibration $> 15\%$}: Flag for review.
        Methods that claim high confidence but produce large errors
        are dangerous in deployment and are flagged for manual
        inspection.
  \item \textbf{Compute budget $> 2\times$ declared}: Automatic
        disqualification.
        Methods that exceed their declared compute budget by more than
        a factor of two are disqualified to prevent undeclared
        resource advantages.
\end{itemize}

\paragraph{Open governance.}
All round reports, RunBundles, scoring code, and failure taxonomies
are published openly.
Disputes are resolved through a structured appeals process with
independent reviewers.
The goal is to build the same kind of institutional trust that CASP
has built in protein structure prediction: trust not in any single
result, but in the \emph{process} that produces results.

\paragraph{Failure taxonomies.}
After each LIP-Arena round, a failure taxonomy is published that
categorizes observed failure modes by type (operator mismatch, solver
limitation, calibration failure, out-of-distribution, compute), by
severity (cosmetic, significant, catastrophic), and by frequency
(isolated, systematic).
These taxonomies serve as a living record of the field's current
limitations and guide both research priorities and future challenge
design.
